\section{Background}

\subsection{Capacity-Constrained Patch Scheduling}

Enterprise vulnerability remediation operates under a hard capacity
constraint: a finite team can patch a finite number of (host, CVE) pairs
per maintenance window. CISA Binding Operational Directive 22-01~\cite{cisa2021bod2201}
mandates 15-day remediation of Known Exploited Vulnerabilities for federal
agencies; NIST SP~800-40 Rev.~4~\cite{nist2022sp80040r4} recommends risk-based
patch management with documented capacity allocation. In practice the
capacity $K$ ---the number of (host, CVE) pairs remediated per window---
is bounded by patching personnel, change-window length, and rollback risk
budgets, and varies by an order of magnitude across organisations.

\subsection{CVE Arrival as Backlog Replenishment}

New CVEs are disclosed continuously: the National Vulnerability Database
publishes new entries daily, and individual organisations see their backlog
grow as scanners surface previously unindexed exposures. The arrival rate
of new \textit{actionable} CVEs per maintenance window ---denoted $\lambda$
in this paper--- depends on fleet size, scanner coverage, and the upstream
CVE disclosure rate. The 2026 DBIR~\cite{verizon2026dbir} reports vulnerability
exploitation as the leading initial-access vector, underscoring the
operational importance of the rate at which actionable disclosures
accumulate.

\subsection{The Capacity-Versus-Arrival Race}

The interaction $(K, \lambda)$ defines a regime: when $K$ greatly exceeds
$\lambda$, the backlog is drained faster than it is replenished, and the
high-EPSS tail of the residual backlog shrinks over time. When $\lambda$
matches or exceeds $K$, the backlog is in steady state and the EPSS
distribution of unpatched pairs remains approximately stationary. Prior temporal-stability work~\cite{paper5temporal}
demonstrated EPSS-only decay in the single regime $(K=50, \lambda=3)$ but
did not characterise the boundary between drain and steady state. This
paper sweeps both axes to map that boundary.

\subsection{EPSS Score Dynamics}

EPSS scores~\cite{jacobs2021epss} are a per-CVE rolling 30-day exploit
probability. Ravalico et al.~\cite{ravalico2025epsstemporal} characterised
per-CVE EPSS drift; their study is complementary to ours --- they measure
the per-CVE score random walk; we measure how prioritisation effectiveness
degrades through compositional change of the residual backlog. The two
effects are independent: the compositional collapse documented here would
occur even with perfectly stable per-CVE EPSS values.
